\chapter{Alternativas de Solución para la Prevención de Accidentes e Incendios}
\label{sec.cap7:alternativas_solucion}

La protección efectiva de los trabajadores y de los activos productivos no puede descansar de forma exclusiva en la
respuesta reactiva ante siniestros ya declarados. En la práctica de la higiene y seguridad en el trabajo, el diseño
ingenieril debe priorizar la erradicación de las condiciones inseguras y de las fuentes potenciales de ignición,
actuando antes de que se complete la cadena causal del accidente. En este capítulo se presentan y describen las
principales alternativas de solución técnicas, organizativas y estructurales orientadas a prevenir focos ígneos y
mitigar sus consecuencias.

\section{Criterios de Clasificación y Jerarquía de Control de Riesgos}
\label{sec.cap7:jerarquia_control}

\subsection{Tipos de Protección contra Incendios}
\begin{itemize}
    \item Protección Preventiva:  Abarca el control estricto de combustibles y comburentes, la eliminación de fuentes de calor, el
          mantenimiento predictivo de instalaciones eléctricas, el orden y limpieza, y los sistemas automáticos de
          detección precoz.
    \item Protección Pasiva o Estructural: Actúa sobre la propagación y la estabilidad del edificio una vez iniciado el
          fuego, sin requerir activación mecánica ni intervención humana. Su propósito es confinar el siniestro, evitar
          el colapso estructural mediante materiales con resistencia al fuego certificada ($F$), impedir el paso de
          humos calientes y garantizar vías de escape seguras.
    \item Protección Activa: Involucra los medios técnicos e instalaciones diseñados para combatir, controlar y
          extinguir el fuego declarado (extintores portátiles, redes fijas de hidrantes, rociadores y sistemas de gas),
          en conjunto con la organización humana de respuesta para la intervención y evacuación.
\end{itemize}

\section{Prevención por Desarticulación de los Factores del Fuego}
\label{sec.cap7:factores_fuego}

Conforme al modelo del tetraedro analizado en el Capítulo \ref{sec.cap2:fenomeno_fisico}, la prevención radica en
impedir la coincidencia temporal y espacial de combustible, comburente y calor con energía de activación suficiente.

\subsection{Medidas sobre el Combustible}

La reducción o supresión de la masa combustible presente en los puestos de trabajo constituye la medida preventiva más
contundente:

\begin{itemize}
    \item Disgregación de elementos: El almacenamiento continuo de grandes volúmenes de materias primas sólidas
          favorece la acumulación interna de calor y la propagación rápida del fuego. Se implementa la técnica de
          disgregación en lotes separadas por pasillos cortafuego de al menos dos metros de ancho, limitando la
          carga térmica por sector e impidiendo la propagación por contacto directo o radiación superficial.
    \item Sustitución de Aceites Minerales: En transformadores de distribución y potencia, el aceite mineral
          hidrocarburado convencional posee un punto de inflamación relativamente bajo. Su reemplazo por aceites
          naturales o sintéticos eleva el punto de combustión por encima de los $300\Celsius$, otorgando propiedades
          que impiden explosiones por arcos internos y reducen distancias de seguridad.
    \item Conductores Libres de Halógenos: El PVC convencional emite densas columnas de hollín y gas al arder. La
          alternativa adoptada es el uso exclusivo de conductores con aislaciones libres de halógenos y de baja emisión
          de humos opacos, los cuales desprenden vapor de agua y retardan la propagación de la llama.
    \item Sustitución de Solventes Inflamables: En tareas de limpieza y desengrase de circuitos impresos y componentes, se
          reemplazan solventes de bajo punto de inflamación (como alcohol isopropílico o acetona) por desengrasantes
          acuosos o por solventes, que carecen de punto de inflamación y anulan la formación de mezclas inflamables.
    \item Almacenamiento Seguro de Sustancias Químicas: Conforme a las directrices de la SRT, los líquidos inflamables se
          almacenan en armarios de seguridad ignífugos homologados, en recipientes metálicos herméticos debidamente
          rotulados, bajo ventilación adecuada y estrictamente distanciados de fuentes calóricas.
\end{itemize}

\subsection{Medidas sobre el Comburente}

En recintos cerrados donde no resulta viable suprimir los combustibles, se actúa sobre la concentración de oxígeno:

\begin{itemize}
    \item Inertización por Manto de Nitrógeno: En tanques de almacenamiento de solventes y reactores de síntesis, se
          inyecta nitrógeno gaseoso comprimido ($N_2$) para mantener una atmósfera inerte en el espacio de cabeza del
          recipiente. Al sostener la concentración de oxígeno por debajo de la Concentración Límite de Oxígeno (LOC $<
          9\percent$), la combustión resulta físicamente imposible aun ante chispas electrostáticas.
    \item Recintos Estancos y Válvulas de Alivio: Los sistemas de almacenamiento de gases y combustibles volátiles se
          diseñan con sellos mecánicos estancos y arrestallamas en líneas de venteo, evitando el contacto con aire libre.
\end{itemize}

\subsection{Medidas sobre el Calor y las Fuentes de Ignición}

El control térmico y eléctrico busca erradicar los focos de encendido en maquinaria e instalaciones:

\begin{itemize}
    \item Mantenimiento Predictivo por Termografía Infrarroja: La resistencia de constricción en bornes y barras de
          distribución flojas genera sobrecalentamientos locales. Se
          ejecutan relevamientos termográficos periódicos bajo carga nominal, detectando gradientes anómalos ($\Delta T \ge
          10\Celsius$) para programar reajustes.
    \item Detección Precoz de Falla de Arco: Para mitigar arcos eléctricos en serie originados
          en conductores deteriorados o fisurados que escapan a la sensibilidad de interruptores termomagnéticos, se
          instalan dispositivos de detección de defecto de arco. Estos equipos analizan las
          firmas de alta frecuencia en la corriente y disparan la apertura del circuito.
    \item Malla de Puesta a Tierra y Descargadores de Sobretensión: Se instala una red con baja resistencia de puesta a
          tierra vinculada a descargadores de sobretensión (DPS).  Esto previene perforaciones dieléctricas por
          descargas atmosféricas y disipa cargas electrostáticas acumuladas.
    \item Mantenimiento de Partes Mecánicas Móviles: La fricción en motores, cintas transportadoras y compresores eleva
          la temperatura superficial por encima del punto de ignición del polvo circundante. Se aplican programas de
          lubricación sistemática y análisis de vibraciones mecánicas.
\end{itemize}

\section{Alternativas Basadas en Controles Administrativos y de Procedimiento}
\label{sec.cap7:administrativos}

Cuando las barreras de ingeniería no eliminan el riesgo residual, los controles procedimentales aseguran una disciplina
estricta en las intervenciones operativas:

\subsection{Procedimiento de Bloqueo y Etiquetado de Energías Peligrosas}

Para intervenir tableros y maquinaria electromecánica sin peligro de descargas de arco o cortocircuitos accidentales, se
aplica el protocolo estandarizado de cinco pasos:
\begin{enumerate}
    \item Notificación formal a los operadores del área comprometida.
    \item Desconexión física y apertura visible de los seccionadores de potencia aguas arriba.
    \item Enclavamiento mecánico mediante cerrojos portacandados con llave individual en poder de cada técnico.
    \item Disipación de energías residuales (descarga de bancos capacitivos y despresurización de líneas neumáticas).
    \item Verificación de ausencia de tensión con detectores bipolares calibrados antes del inicio de los trabajos.
\end{enumerate}

\subsection{Orden, Limpieza y Gestión de Cargas Combustibles (Metodología 5S)}

La falta de orden en talleres y almacenes incrementa la carga de fuego de manera inadvertida. De acuerdo con las
recomendaciones de la SRT, la metodología 5S erradica el acopio no autorizado de embalajes, maderas y plásticos en zonas
de trabajo. Asimismo, garantiza pasillos de evacuación despejados y una zona libre mínima de un metro de radio alrededor
de tableros eléctricos, interruptores generales y nichos de hidrantes.

\section{Medidas de Protección Pasiva y Sectorización Estructural}
\label{sec.cap7:proteccion_pasiva}

La protección pasiva limita la propagación del incendio y evita el colapso del establecimiento mediante soluciones
constructivas:

\begin{itemize}
    \item Sectorización de Incendios: División del edificio en recintos independientes delimitados por muros cortafuego con
          resistencia al fuego certificada ($F30$, $F60$, $F90$, $F120$ o $F180$) conforme al cálculo de carga de fuego
          establecido en el Anexo VII del Decreto 351/79.
    \item Puertas Cortafuego de Doble Contacto: Instalación de puertas resistentes al fuego con cierre automático y sellos
          perimetrales intumescentes. Estos sellos expanden su volumen ante temperaturas elevadas, impidiendo el pasaje de
          humos fríos y gases tóxicos calientes hacia las cajas de escalera.
    \item Sellado de Pasajes de Instalaciones: Relleno de pases de bandejas portacables y tuberías a través de tabiques
          con masillas intumescentes, almohadillas cortafuego y collares elásticos que obturan la abertura si las cubiertas
          de los cables se consumen.
    \item Vías de Escape y Cajas de Escalera Presurizadas: Mantenimiento de pasillos de evacuación anchos y libres de
          combustibles, complementados en edificios en altura con sistemas mecánicos de inyección de aire que mantienen
          una sobrepresión positiva ($50\,\mathrm{Pa}$) en la caja de escaleras, asegurando un ambiente limpio de humo
          para el descenso de los ocupantes.
\end{itemize}

\section{Organización Humana de Respuesta y Equipos de Intervención}
\label{sec.cap7:organizacion_humana}

La existencia de recursos materiales resulta ineficaz si el factor humano carece de una estructura organizativa clara y
disciplinada para la respuesta ante contingencias. La doctrina técnica de prevención estructura los medios humanos en
dos niveles escalonados y un órgano de conducción:

\subsection{Equipo de Primera Intervención (E.P.I.)}

El Equipo de Primera Intervención constituye la fuerza de choque inmediata en el sector o taller donde se origina el
foco. Se conforma con trabajadores habituales de cada área seleccionados por su conocimiento del proceso y aptitud:

\begin{itemize}
    \item Misión de Advertencia: Al advertir un principio de fuego, su primera obligación consiste en accionar el
          pulsador de alarma más próximo, comunicar la emergencia a la central de control por la línea interna asignada y
          alertar al servicio de bomberos externo.
    \item Misión de Intervención: Atacar de inmediato el fuego en su fase incipiente empleando los medios ligeros del área
          (extintores portátiles o mangueras de la Red Interior de Agua - R.I.A.), priorizando la propia seguridad y la de
          sus compañeros. Si el fuego escapa a las capacidades del extintor, repliegan ordenadamente asegurando el cierre
          de puertas para confinar el siniestro.
\end{itemize}

\subsection{Equipo de Segunda Intervención (E.S.I.)}

El Equipo de Segunda Intervención conforma la brigada contra incendios propiamente dicha de la empresa. Está compuesto por
personal con un adiestramiento superior, continuo y periódico, capacitado para intervenir en cualquier sector de la
planta:

\begin{itemize}
    \item Concentración Rápida: Ante la alarma general, sus miembros acuden al punto de reunión prefijado para equiparse
          con elementos de protección personal pesados y equipos de respiración autónoma (ERA).
    \item Ataque con Medios Pesados: Operan líneas de mangueras de $45\mm$ y $63.5\mm$ desde hidrantes, equipos generadores
          de espuma, realizan el corte selectivo de suministros energéticos y ejecutan maniobras de rescate de operarios
          atrapados.
    \item Coordinación con Bomberos Oficiales: Al arribo de los cuerpos de bomberos públicos, el jefe del E.S.I. informa la
          ubicación exacta del foco, las materias primas involucradas, los accesos seguros y el estado del recuento de
          personal.
\end{itemize}

\subsection{Comité de Emergencia y Plan de Evacuación}

La dirección de la contingencia recae en el Comité de Emergencia encabezado por el Jefe de Emergencia. Su función radica
en evaluar la magnitud del incidente, ordenar la evacuación total o sectorial de las instalaciones, supervisar la
apertura de accesos para vehículos de socorro y verificar el recuento del personal en el Punto de Encuentro Seguro
establecido en el exterior del establecimiento.

